(1) 在极坐标中角动量 \(\ell=mr^2\dot\phi\) 守恒，径向有效势为 \(U_{\rm eff}=-K/r+(\ell^2/m+\alpha)/(2r^2)\)。令 \(u=1/r\)，径向力 \(F_r=-K/r^2+\alpha/r^3\)，Binet方程 \(u''+u=-mF_r(1/u)/(\ell^2u^2)\) 给 \(u''+[1+m\alpha/\ell^2]u=mK/\ell^2\)。
(2) 记 \(\beta=\sqrt{1+m\alpha/\ell^2}>0\)，平衡倒半径 \(u_0=mK/(\ell^2+m\alpha)>0\)。于是 \(u(\phi)=u_0[1+e\cos(\beta(\phi-\phi_0))]\)，\(0\le e<1\) 保证分母始终正、径向轨道有界，\(\phi_0\) 决定一处近心点方向。
(3) 余弦在自变量变化 \(2\pi\) 后重复，因此两次近心点方位差为 \(\Delta\phi_{\rm aps}=2\pi/\beta\)。相对开普勒情形的近心点进动量为 \(\Delta\varpi=2\pi(1/\beta-1)\)；当 \(\alpha>0\) 时 \(\beta>1\)，近心点后退，\(\alpha<0\) 且稳定条件成立时则进动超前。
(4) 有效势的圆轨道满足 \(Kr_c=\ell^2/m+\alpha\)，所以 \(r_c=(\ell^2+m\alpha)/(mK)>0\)。二阶导 \(U_{\rm eff}''(r_c)=-2K/r_c^3+3(\ell^2/m+\alpha)/r_c^4=K/r_c^3>0\)，故稳定。若 \(|\alpha|m\ll\ell^2\)，展开 \((1+m\alpha/\ell^2)^{-1/2}\simeq1-m\alpha/(2\ell^2)\)，得 \(\Delta\varpi\simeq-\pi m\alpha/\ell^2\)。
